\scanpage{143}
\[
(8)\qquad (\sigma^\nu.u).t_\varepsilon=(-1)^\nu(t+u)_\varepsilon
\]
\margin{(NB $t_\varepsilon=(t,\varepsilon)$\quad $\varepsilon\in\{\pm1\}$)}
Cette opération passe au quotient en une opération de $D_{\mathbb R}$ sur $\deleted{\unsure{\mathbb R}}\mathbb U_\varepsilon$ et $\overline{\mathrm u}$ de $D_{\mathbb U}$ sur $\mathbb U_\varepsilon$, cette dernière est donnée par
\[
(9)\qquad (\sigma^\nu.\theta)z_\varepsilon=(z\theta)^{(-1)^\nu}=\left\{\begin{array}{ll}
z\theta&\text{si }\nu\equiv0\quad(2)\\
z^{-1}\theta^{-1}&\text{si }\nu\equiv1\quad(2).
\end{array}\right.
\]
Ainsi le groupe $D_{\mathbb R}\times D_{\mathbb R}$ opère sur
\[
\widetilde{\partial X}\overset{\mathrm{dfn}}{=}\widetilde{\mathbb U}_{-1}\amalg\widetilde{\mathbb U}_1\simeq\unsure{\widetilde{\mathbb R}_{-1}\amalg\widetilde{\mathbb R}_1},
\]
via des actions sur les deux sommandes, et cette action passe au quotient en une action de $D_{\mathbb U}\times D_{\mathbb U}$ sur $\deleted{\unsure{\partial X}}\partial X=\mathbb U_{-1}\amalg\mathbb U_1$. Nous allons prolonger cette action de $D_{\mathbb R}\times D_{\mathbb R}$ sur $\widetilde{\partial X}$ en une action sur
\[
(10)\qquad \widetilde X\overset{\mathrm{dfn}}{=}\mathbb R\times\unsure{\mathbb I|\mathbb J}
\]
qui revêt $X$ par
\[
(11)\qquad (t,s)\longmapsto(\exp2i\pi t,s)\ ;
\]
\deleted{Pour ceci} Plus exactement, c'est le sous-groupe
\[
(12)\qquad D_{\mathbb R\times\mathbb R}\overset{\mathrm{dfn}}{=}\underset{(1/2\text{ direct})}{\{1,\sigma\}.(\mathbb R\times\mathbb R)}\hookrightarrow D_{\mathbb R}\times D_{\mathbb R}
\]
\[
-\sigma^\nu.(\unsure{u_{-1}|t_{-1}},\unsure{u_1|t_1})\mathbin{\mapstochar\!\relbar\!\relbar}(\sigma^\nu.\unsure{u_{-1}|t_{-1}},\sigma^\nu.\unsure{u_{+1}|t_{+1}})
\]
qu'on va faire opérer, de façon à ce que ces opérations passent au quotient en une opération de $D_{\mathbb R\times\mathbb R}$ sur $X$ lui-$\overline{\mathrm m}$. On
\margin{143}
